% \section{Wind-Conditioned Source Apportionment}
% \label{sec:wind_conditioned_source_apportionment}

% Section~\ref{sec:problem_setup} formulated source apportionment as a
% lagged inverse problem over known source inventories. We now describe how
% the lagged source-response matrix is constructed in practice from sparse
% wind observations, source inventories, a physically constrained transport
% operator, and background correction. The output of this section is a
% projected source-response matrix and a nonnegative estimate of source
% contributions. Section~\ref{sec:identifiability_theory} then analyzes
% when these estimates are identifiable.

% \subsection{Method Overview}

% The proposed apportionment model has four stages. First, sparse wind
% measurements are interpolated to the spatial grid. Second, the gridded
% wind field is used to construct an open-boundary lagged transport
% operator that maps emissions released at earlier times to concentrations
% at later sensor times. Third, the transported source inventories are
% sampled at the sensor locations to form source fingerprints. Fourth,
% background and temporal components are removed before estimating
% nonnegative source activity multipliers.

% The central object is the projected lagged source-response matrix
% \begin{equation}
% \widetilde H
% :=
% P_Q^\perp H_{1:T}^{\mathrm{lag}},
% \qquad
% \widetilde{\mathbf Y}
% :=
% P_Q^\perp \mathbf Y,
% \label{eq:projected_HY_def}
% \end{equation}
% where \(H_{1:T}^{\mathrm{lag}}\) encodes the wind-conditioned response of
% each source group, and \(P_Q^\perp\) removes low-dimensional background
% and temporal components. The source contribution vector is then estimated
% by solving a constrained inverse problem of the form
% \begin{equation}
% \widehat{\boldsymbol\theta}
% =
% \arg\min_{\boldsymbol\theta\ge 0}
% \left\|
% \widetilde{\mathbf Y}
% -
% \widetilde H\boldsymbol\theta
% \right\|_2^2
% +
% \lambda R(\boldsymbol\theta).
% \label{eq:projected_nnls_estimator_overview}
% \end{equation}

% \subsection{Coordinate-Aware Wind Field Estimation}

% Wind is typically observed at a small number of monitoring locations and
% must be estimated on the full concentration grid. Let
% \[
% X_{\mathrm{wind}}
% =
% \{\mathbf z_1,\ldots,\mathbf z_p\}
% \]
% denote the locations at which wind measurements are available, and let
% \(\mathbf w_t^{\mathrm{obs}}(\mathbf z_i)\in\mathbb R^2\) denote the
% observed horizontal wind vector at time \(t\). We estimate a gridded wind
% field
% \[
% \widehat{\mathbf w}_t(\mathbf x),
% \qquad
% \mathbf x\in\Omega,
% \]
% using a coordinate-aware interpolation model
% \begin{equation}
% \widehat{\mathbf w}_{1:T}
% =
% f_\phi
% \left(
% \mathbf w_{1:T}^{\mathrm{obs}},
% X_{\mathrm{wind}},
% X_{\mathrm{grid}}
% \right).
% \label{eq:wind_interpolation_model}
% \end{equation}

% The interpolation model may be implemented using kriging, inverse-distance
% weighting, a graph interpolation model, or a local transformer-style
% model over nearby wind observations. Since the measurements are hourly,
% we do not require a long-horizon temporal sequence model. Instead, the
% model may use a short temporal patch around each hour together with
% spatial coordinates. This is important because wind interpolation is not
% determined by locality alone: terrain, urban geometry, and coordinates
% can affect the wind field even when nearby observations are sparse.

% In practice, we use a residual parameterization around a physically
% reasonable baseline:
% \begin{equation}
% \widehat{\mathbf w}_t(\mathbf x)
% =
% \mathbf w_t^{0}(\mathbf x)
% +
% \Delta_\phi(\mathbf x,t),
% \label{eq:wind_residual_parameterization}
% \end{equation}
% where \(\mathbf w_t^{0}\) is a baseline interpolation, such as kriging or
% reanalysis-guided interpolation, and \(\Delta_\phi\) is a bounded learned
% correction. The wind model is first trained using wind measurements alone:
% \begin{equation}
% \begin{aligned}
% \mathcal L_{\mathrm{wind}}
% =
% &\sum_{t}
% \sum_{\mathbf z_i\in X_{\mathrm{wind}}}
% \left\|
% \widehat{\mathbf w}_t(\mathbf z_i)
% -
% \mathbf w_t^{\mathrm{obs}}(\mathbf z_i)
% \right\|_2^2
% \\
% &+
% \lambda_{\mathrm{smooth}}
% \mathcal R_{\mathrm{smooth}}(\widehat{\mathbf w})
% +
% \lambda_{\Delta}
% \|\Delta_\phi\|_2^2 .
% \end{aligned}
% \label{eq:wind_pretraining_loss}
% \end{equation}
% We use a spatial smoothness regularizer on the gridded wind field:
% \begin{equation}
% \mathcal R_{\mathrm{smooth}}(\widehat{\mathbf w})
% =
% \sum_t
% \sum_{(g,g')\in \mathcal E_{\mathrm{grid}}}
% \left\|
% \widehat{\mathbf w}_t(\mathbf x_g)
% -
% \widehat{\mathbf w}_t(\mathbf x_{g'})
% \right\|_2^2,
% \label{eq:wind_spatial_smoothness}
% \end{equation}
% where \(\mathcal E_{\mathrm{grid}}\) denotes neighboring grid-cell pairs.
% This term discourages physically implausible cell-to-cell oscillations in
% the interpolated wind field while still allowing spatial variation driven
% by the observed wind measurements and coordinates.

% The final term \(\|\Delta_\phi\|_2^2\) penalizes deviation from the
% baseline wind interpolation in the residual parameterization
% \(\widehat{\mathbf w}_t=\mathbf w_t^0+\Delta_\phi\). This pretraining
% step prevents the wind field from becoming an unconstrained latent
% variable that adapts only to fit pollution sensors.

% \subsection{Open-Boundary Lagged Plume Response}

% We now define the transport operator used to map source emissions to
% concentration fields. Let \(G_{t,\tau}^{\partial}\) denote the
% open-boundary transport response from release time \(\tau\) to sensor
% time \(t\), where \(\tau\le t\). The superscript \(\partial\) indicates
% that the computational domain is open: pollutant mass that is advected or
% dispersed outside \(\Omega\) leaves the accounting and is not reflected,
% wrapped around, or redistributed inside the domain.

% For a nonnegative emission vector \(\mathbf z\in\mathbb R^q_+\), the
% transport operator satisfies
% \begin{equation}
% G_{t,\tau}^{\partial}\mathbf z \ge 0,
% \qquad
% \mathbf 1^\top
% G_{t,\tau}^{\partial}\mathbf z
% \le
% \mathbf 1^\top \mathbf z.
% \label{eq:open_boundary_mass_nonincrease}
% \end{equation}
% Thus, the operator preserves nonnegativity and is mass non-increasing
% over the modeled domain.

% We instantiate \(G_{t,\tau}^{\partial}\) using a differentiable
% lagged puff/plume approximation. Consider a unit emission released from
% source-grid cell \(i\) at location \(\mathbf r_i\) and time \(\tau\).
% The puff center is initialized as
% \begin{equation}
% \mathbf z_i(\tau)=\mathbf r_i.
% \end{equation}
% For intermediate times \(a\in[\tau,t]\), the puff center is advected by
% the interpolated wind field:
% \begin{equation}
% \frac{d\mathbf z_i(a)}{da}
% =
% \widehat{\mathbf w}_a(\mathbf z_i(a)).
% \label{eq:puff_center_advection}
% \end{equation}
% In a discrete implementation with substep size \(\delta a\), this becomes
% \begin{equation}
% \mathbf z_i^{r+1}
% =
% \mathbf z_i^r
% +
% \delta a\,
% \widehat{\mathbf w}_{a_r}(\mathbf z_i^r).
% \label{eq:discrete_puff_center_advection}
% \end{equation}
% If the puff center exits \(\Omega\), the corresponding in-domain mass is
% removed. If the puff remains inside \(\Omega\), its contribution at time
% \(t\) is spread using a Gaussian dispersion kernel
% % \begin{equation}
% % K_\psi
% % \left(
% % \mathbf x;
% % \mathbf z_i(t),
% % \Sigma_i(t,\tau)
% % \right)
% % =
% % \frac{1}{2\pi |\Sigma_i(t,\tau)|^{1/2}}
% % \exp\left(
% % -\frac{1}{2}
% % (\mathbf x-\mathbf z_i(t))^\top
% % \Sigma_i(t,\tau)^{-1}
% % (\mathbf x-\mathbf z_i(t))
% % \right),
% % \label{eq:gaussian_puff_kernel}
% % \end{equation}
% \begin{equation}
% K_\psi(\mathbf x;\mathbf z_i,\Sigma_i)
% =
% \frac{
% \exp\!\left(
% -\frac{1}{2}\|\mathbf x-\mathbf z_i\|_{\Sigma_i^{-1}}^2
% \right)
% }{
% 2\pi |\Sigma_i|^{1/2}
% }.
% \label{eq:gaussian_puff_kernel}
% \end{equation}
% \[
% \|\mathbf a\|_{\Sigma^{-1}}^2 := \mathbf a^\top \Sigma^{-1}\mathbf a .
% \]
% where \(\psi\) denotes dispersion parameters and
% \(\Sigma_i(t,\tau)\) controls downwind and crosswind spread. A simple
% choice is
% \begin{equation}
% \Sigma_i(t,\tau)
% =
% R_i(t,\tau)
% \begin{bmatrix}
% \sigma_{\parallel}^2(t-\tau) & 0\\
% 0 & \sigma_{\perp}^2(t-\tau)
% \end{bmatrix}
% R_i(t,\tau)^\top,
% \label{eq:anisotropic_dispersion_covariance}
% \end{equation}
% where \(R_i(t,\tau)\) aligns the covariance with the mean wind direction
% along the puff trajectory. The dispersion scales
% \(\sigma_{\parallel}\) and \(\sigma_{\perp}\) may be fixed from plume
% physics or calibrated within physically plausible bounds.

% Let \(\{\mathbf x_g\}_{g=1}^n\) denote the concentration-grid cell
% centers. The response of the grid field to a unit release from source
% cell \(i\) is
% \begin{equation}
% \left[
% G_{t,\tau}^{\partial}
% \left(
% \widehat{\mathbf w}_{\tau:t};\psi
% \right)
% \mathbf e_i
% \right]_g
% =
% \chi_i(t,\tau)
% K_\psi
% \left(
% \mathbf x_g;
% \mathbf z_i(t),
% \Sigma_i(t,\tau)
% \right)
% \Delta V_g,
% \label{eq:puff_operator_column}
% \end{equation}
% where \(\mathbf e_i\) is the \(i\)-th source basis vector,
% \(\Delta V_g\) is the grid-cell area or volume factor, and
% \(\chi_i(t,\tau)\in\{0,1\}\) indicates whether the puff remains active in
% the open domain. Kernel mass outside \(\Omega\) is not renormalized, so
% dispersion across the boundary naturally removes mass from the modeled
% domain.

% This construction implements a
% mass-accounting transport response. Faster winds move the puff center
% farther during a fixed time interval, while dispersion controls how
% broadly the emitted mass is distributed across grid cells. The sensor
% signal then depends on whether the advected and dispersed plume overlaps
% the monitoring locations at later times.

% \subsection{Lagged Inventory-to-Sensor Response}

% Let \(S\in\mathbb R^{q\times K}\) be the source-inventory matrix, where
% the \(k\)-th column \(\mathbf s_k\) represents the spatial pattern of
% source group \(k\). For a fixed source activity vector
% \(\boldsymbol\theta\in\mathbb R^K_+\), emissions released during an
% interval are proportional to \(S\boldsymbol\theta\).

% At sensor time \(t\), emissions released at earlier times
% \(t-\ell\) may still contribute to the observed concentration. Let
% \(\mathcal L_t=\{0,\ldots,\min(L,t-1)\}\) be the available lag set. The
% lagged source-response matrix is
% \begin{equation}
% H_t^{\mathrm{lag}}
% =
% \sum_{\ell\in\mathcal L_t}
% \mathcal O
% G_{t,t-\ell}^{\partial}
% \left(
% \widehat{\mathbf w}_{t-\ell:t};\psi
% \right)
% S
% \in\mathbb R^{m\times K}.
% \label{eq:method_lagged_response_matrix}
% \end{equation}
% Stacking over time gives
% \begin{equation}
% H_{1:T}^{\mathrm{lag}}
% =
% \begin{bmatrix}
% H_1^{\mathrm{lag}}\\
% H_2^{\mathrm{lag}}\\
% \vdots\\
% H_T^{\mathrm{lag}}
% \end{bmatrix}
% \in
% \mathbb R^{mT\times K}.
% \label{eq:method_stacked_lagged_response_matrix}
% \end{equation}

% The \(k\)-th column of \(H_{1:T}^{\mathrm{lag}}\) is the finite-horizon
% sensor fingerprint of source group \(k\) under the estimated wind field,
% open-boundary transport, and chosen lag window. This fingerprint is not
% only spatial but spatio-temporal: it records when and where emissions
% from a source group become visible at the sparse sensor network.

% \subsection{Background and Temporal Correction}

% Observed pollution includes regional background, time-of-day variation,
% seasonal trends, low-frequency drift, and sensor offsets that are not
% explained by the local source inventory. We represent these components
% using a low-dimensional background-correction basis
% \begin{equation}
% Q\in\mathbb R^{mT\times r}.
% \end{equation}
% The columns of \(Q\) may include cubic spline basis functions over
% time-of-day, day-level intercepts, regional background trends, and
% sensor-specific offsets.

% The observation model is
% \begin{equation}
% \mathbf Y
% =
% H_{1:T}^{\mathrm{lag}}\boldsymbol\theta
% +
% Q\boldsymbol\gamma
% +
% \mathbf E,
% \label{eq:method_background_model}
% \end{equation}
% where \(\boldsymbol\gamma\in\mathbb R^r\) are background-correction
% coefficients. To separate the local source-response signal from the
% background subspace, we use the projection
% \begin{equation}
% P_Q^\perp
% =
% I
% -
% Q Q^\dagger,
% \label{eq:method_background_projection}
% \end{equation}
% where \(Q^\dagger\) denotes the Moore--Penrose pseudoinverse. Applying
% this projection gives
% \begin{equation}
% \widetilde{\mathbf Y}
% =
% P_Q^\perp \mathbf Y,
% \qquad
% \widetilde H
% =
% P_Q^\perp H_{1:T}^{\mathrm{lag}}.
% \label{eq:method_projected_data_matrix}
% \end{equation}

% The choice of \(Q\) must be deliberately low-dimensional. If the
% background basis is too flexible, it may absorb source-driven variation
% and reduce the identifiable source signal. In our implementation, the
% background correction is therefore restricted to smooth temporal
% components and low-rank sensor offsets rather than arbitrary
% high-frequency per-sensor curves.

% \subsection{Source Activity Estimation}

% Given the projected data and response matrix, we estimate source
% contributions by solving a nonnegative regularized least-squares problem:
% \begin{equation}
% \widehat{\boldsymbol\theta}
% =
% \arg\min_{\boldsymbol\theta\ge 0}
% \left\|
% \widetilde{\mathbf Y}
% -
% \widetilde H\boldsymbol\theta
% \right\|_2^2
% +
% \lambda R(\boldsymbol\theta).
% \label{eq:method_source_estimator}
% \end{equation}
% The regularizer \(R\) can encode weak prior information about source
% activity. In the simplest case,
% \begin{equation}
% R(\boldsymbol\theta)
% =
% \|\boldsymbol\theta\|_2^2,
% \label{eq:ridge_source_regularizer}
% \end{equation}
% which stabilizes estimation when source fingerprints are nearly
% collinear. If prior activity levels \(\boldsymbol\theta_0\) are available,
% we may instead use
% \begin{equation}
% R(\boldsymbol\theta)
% =
% \|\boldsymbol\theta-\boldsymbol\theta_0\|_2^2.
% \label{eq:inventory_prior_regularizer}
% \end{equation}
% Other choices, such as sparsity penalties, can be used when only a small
% number of source groups are expected to dominate the observed signal.

% Equivalently, one may fit source and background coefficients jointly:
% \begin{equation}
% (\widehat{\boldsymbol\theta},\widehat{\boldsymbol\gamma})
% =
% \arg\min_{\boldsymbol\theta\ge 0,\boldsymbol\gamma}
% \left\|
% \mathbf Y
% -
% H_{1:T}^{\mathrm{lag}}\boldsymbol\theta
% -
% Q\boldsymbol\gamma
% \right\|_2^2
% +
% \lambda R(\boldsymbol\theta).
% \label{eq:joint_source_background_fit}
% \end{equation}
% The projected form in \eqref{eq:method_source_estimator} is useful
% because it makes explicit which part of the source-response matrix
% remains after background correction. This projected matrix is the object
% used for the identifiability analysis in the next section.

% \subsection{Constrained End-to-End Refinement}

% The core estimator treats the wind field, transport parameters, source
% inventories, and background basis as fixed when estimating
% \(\boldsymbol\theta\). This is the most interpretable setting. However,
% when wind interpolation or dispersion parameters are uncertain, we allow
% a constrained refinement stage after the initial fit.

% Let \(\phi_0\) denote the pretrained wind-interpolation parameters and
% \(\psi_0\) denote initial plume or dispersion parameters. We refine
% \((\phi,\psi,\boldsymbol\theta,\boldsymbol\gamma)\) by minimizing
% \begin{align}
% \mathcal L_{\mathrm{refine}}
% =
% &
% \left\|
% \mathbf Y
% -
% H_{1:T}^{\mathrm{lag}}(\phi,\psi)\boldsymbol\theta
% -
% Q\boldsymbol\gamma
% \right\|_2^2
% +
% \lambda_\theta R(\boldsymbol\theta)
% \nonumber\\
% &
% +
% \lambda_w
% \left\|
% \widehat{\mathbf w}_{1:T}^{\,\phi}
% -
% \widehat{\mathbf w}_{1:T}^{\,\phi_0}
% \right\|_2^2
% +
% \lambda_\psi
% \|\psi-\psi_0\|_2^2
% \nonumber \\
% &
% +
% \lambda_{\mathrm{smooth}}
% \mathcal R_{\mathrm{smooth}}(\widehat{\mathbf w}^{\,\phi}),
% \label{eq:constrained_refinement_loss}
% \end{align}
% subject to
% \begin{equation}
% \boldsymbol\theta\ge 0,
% \qquad
% \psi\in\Psi_{\mathrm{phys}},
% \qquad
% \left\|
% \widehat{\mathbf w}_{1:T}^{\,\phi}
% -
% \widehat{\mathbf w}_{1:T}^{\,\phi_0}
% \right\|_\infty
% \le
% \epsilon_w.
% \label{eq:refinement_constraints}
% \end{equation}
% Here \(\Psi_{\mathrm{phys}}\) denotes physically plausible bounds on
% dispersion and plume parameters.

% This refinement is intentionally constrained. Without such constraints,
% the wind field and plume parameters can become flexible latent variables
% that fit sparse pollution sensors while weakening the interpretability of
% the inferred source contributions. In experiments, we therefore report
% both the fixed-operator estimate and the constrained-refinement estimate.

% % \subsection{Outputs of the Apportionment Model}

% % The method produces the following quantities:
% % \begin{enumerate}[leftmargin=*]
% %     \item source contribution estimates
% %     \(\widehat{\boldsymbol\theta}\);
% %     \item fitted sensor trajectories
% %     \[
% %     \widehat{\mathbf Y}
% %     =
% %     H_{1:T}^{\mathrm{lag}}\widehat{\boldsymbol\theta}
% %     +
% %     Q\widehat{\boldsymbol\gamma};
% %     \]
% %     \item projected residuals
% %     \[
% %     \widetilde{\mathbf r}
% %     =
% %     \widetilde{\mathbf Y}
% %     -
% %     \widetilde H\widehat{\boldsymbol\theta};
% %     \]
% %     \item projected source fingerprints, given by the columns of
% %     \(\widetilde H\).
% % \end{enumerate}

% The outputs of this approach are not sufficient by themselves to justify source-level claims. A source estimate may fit the sensor data well while still being
% unstable if the corresponding source fingerprints are weak, collinear, or
% absorbed by background correction. The next section therefore studies the
% rank, conditioning, visibility, and coherence of \(\widetilde H\), which
% determine the identifiable resolution of the source apportionment
% problem.


\section{Wind-Conditioned Response Construction and Estimation}
\label{sec:wind_conditioned_source_apportionment}

Section~\ref{sec:problem_setup} defines the projected inverse problem.  This
section describes how the response matrix is constructed in practice and how
source activities are estimated once the response has been formed.

\subsection{Masked Wind Preparation}
\label{subsec:wind_field_estimation}

Government wind direction is meteorological: it records where wind comes from,
whereas transport requires the direction in which pollution moves.  We first
convert direction \(WD\) in degrees and speed \(WS\) to transport vectors,
\begin{equation}
\begin{aligned}
U_x&=-WS\sin(WD\pi/180),\\
V_y&=-WS\cos(WD\pi/180).
\end{aligned}
\label{eq:wind_direction_conversion}
\end{equation}
The observed vectors and their masks are supplied to a fixed-node masked
ImputeFormer model \citep{nie2024imputeformer}.  Training masks additional
observed entries, while inference preserves every observed vector and fills
only missing entries.  The completed station vectors are averaged at each hour
to produce the city-level sequence used by the first New Delhi response path,
\begin{equation}
\widehat{\mathbf w}_t
=\frac{1}{p}\sum_{i=1}^{p}\widehat{\mathbf w}_t(\mathbf z_i).
\label{eq:city_wind_average}
\end{equation}
This fixed-node, city-level model is the implementation used for the study; a
coordinate-aware spatial wind field is a future extension.

Meteorological speed must be converted to grid displacement before puff
advection.  For response timestep \(\Delta t_{\mathrm{s}}\) seconds and cell
dimensions \(\Delta x_{\mathrm{m}},\Delta y_{\mathrm{m}}\), we use
\begin{equation}
v_x^{\mathrm{grid}}(t)=\frac{U_x(t)\Delta t_{\mathrm{s}}}{\Delta x_{\mathrm{m}}},
\qquad
v_y^{\mathrm{grid}}(t)=\frac{V_y(t)\Delta t_{\mathrm{s}}}{\Delta y_{\mathrm{m}}}.
\label{eq:wind_unit_conversion}
\end{equation}
Every response artifact records these scales, the coordinate convention, the
wind units, and the exact converted sequence.  Controlled experiments use the
same interface for constant, single-direction, diurnal, AR(1), and
multi-directional synthetic winds.

\subsection{Open-Boundary Lagged Plume Response}
\label{subsec:plume_response}

The transport response \(G_{t,\tau}^{\partial}\) maps emissions released at
\(\tau\le t\) to the concentration field at time \(t\).  For any nonnegative
emission vector \(\mathbf z\in\mathbb R_+^q\), the open-boundary response obeys
\begin{equation}
G_{t,\tau}^{\partial}\mathbf z\ge 0,
\qquad
\mathbf 1^\top G_{t,\tau}^{\partial}\mathbf z\le \mathbf 1^\top\mathbf z.
\label{eq:mass_nonincreasing}
\end{equation}
Thus the operator preserves nonnegativity and is mass non-increasing inside the
modeled domain.

We instantiate this response with a Gaussian puff approximation.  For a
unit release from source cell \(i\) at location \(\mathbf r_i\) and release time
\(\tau\), the puff center is initialized as \(\mathbf z_i(\tau)=\mathbf r_i\) and
advected by the interpolated wind field,
\begin{equation}
\frac{d\mathbf z_i(a)}{da}=\widehat{\mathbf w}_a(\mathbf z_i(a)),
\qquad a\in[\tau,t].
\label{eq:puff_ode}
\end{equation}
Source cells have integer centers and domain extents
\([-0.5,N_x-0.5]\times[-0.5,N_y-0.5]\).  With substep size \(\delta a\), the
trajectory is discretized as
\begin{equation}
\mathbf z_i^{r+1}=\mathbf z_i^r+\delta a\,\widehat{\mathbf w}_{a_r}(\mathbf z_i^r).
\label{eq:puff_discrete}
\end{equation}
The final substep is shortened so that the trajectory lands exactly at the
requested observation time.  If the center exits \(\Omega\), the whole
remaining puff is removed and never reflected, wrapped, clamped, or reinserted.
Otherwise its
contribution is spread with an anisotropic Gaussian kernel
\begin{equation}
K_\psi(\mathbf x;\mathbf z_i,\Sigma_i)
=\frac{\exp\left(-\frac12\|\mathbf x-\mathbf z_i\|_{\Sigma_i^{-1}}^2\right)}
{2\pi |\Sigma_i|^{1/2}},
\label{eq:gaussian_kernel}
\end{equation}
where \(\psi\) denotes dispersion parameters.  A simple covariance model is
\begin{equation}
\Sigma_i(t,\tau)
= R_i(t,\tau)
\begin{bmatrix}
\sigma_\parallel^2(t-\tau) & 0\\
0 & \sigma_\perp^2(t-\tau)
\end{bmatrix}
R_i(t,\tau)^\top,
\label{eq:dispersion_covariance}
\end{equation}
with \(R_i\) aligned to the mean wind direction along the trajectory and a
fixed fallback orientation when mean wind is numerically zero.  The age in this
covariance is lower-bounded by a positive minimum dispersion time, so
same-time releases remain finite.

Response-matrix entries are computed by evaluating the kernel directly at
sensor coordinates:
\begin{equation}
\left[O G_{t,\tau}^{\partial}
(\widehat{\mathbf w}_{\tau:t};\psi)\mathbf e_i\right]_s
=\chi_i(t,\tau)K_\psi(\mathbf x_s;\mathbf z_i(t),\Sigma_i(t,\tau)).
\label{eq:sensor_response}
\end{equation}
Grid quadrature is used separately for kernel retained/dropped-mass diagnostics.
Kernel boundary or truncation loss and whole-release exit loss are recorded
separately and are never renormalized.  Summing repeated sensor observations is
not a mass estimate.

Paper-facing response matrices use a zero-source, zero-initial-state baseline.
Kriged initial conditions belong only to auxiliary PDE mismatch experiments;
their matched zero-source trajectory is subtracted before attribution.  The
operator-matched experiments use the open-boundary puff response, while the
edge-hold PDE and perturbed puff operators are reserved for explicitly labeled
response-error tests.

The computational roadmap implements this response in PyTorch.  Tensor
operations retained in the graph can be differentiated, but discrete exit
events and reporting logic are not assigned gradients.

\subsection{Background Construction and Projection}
\label{subsec:background_construction}

The normal background basis may contain a global constant, centered and scaled
elapsed-time polynomials, local-clock daily sine/cosine harmonics,
reference-coded day effects, standardized regional sensor-coordinate trends,
reference-coded sensor offsets, and a row-aligned user basis.  It depends only
on timestamps, day labels, sensor identity, and sensor coordinates.  Its
paper-facing effective rank is capped at eight.  Source-derived columns are
allowed only in labeled stress tests.

The primary background specification is declared before source fitting and is
not selected by inspecting \(Y\) or by optimizing source recovery.  For the New
Delhi study it is the rank-four basis consisting of a global constant, centered
linear trend, and first daily sine and cosine.  Alternative source-independent
bases are predeclared sensitivity analyses; a source-like basis is only a
labeled stress test.  None is data-selected as a replacement authoritative
model.  Each comparison reports the resulting visibility, absorption, rank,
conditioning, and ambiguity components.

Let \(Q=U\Sigma V^\top\) be a thin singular value decomposition and retain
columns \(U_r\) above the recorded rank tolerance.  Projection is applied
implicitly,
\begin{equation}
P_Q^\perp X=X-U_r(U_r^\top X),
\label{eq:implicit_background_projection}
\end{equation}
without constructing an \(N\times N\) matrix.  Metadata record requested and
effective rank, discarded/dependent columns, tolerance, and exact row order.
The response, observation, and background row records must match exactly.

\subsection{Response Matrix and Source Activity Fit}
\label{subsec:response_and_fit}

Given the estimated wind field, source inventory, and temporal activity basis,
the lagged source--basis response is constructed as in
Section~\ref{sec:problem_setup}.  Its \((k,b)\)-th column is
\begin{equation}
\mathbf h_{t,kb}^{\mathrm{lag}}
=\sum_{\ell\in\mathcal L_t}\phi_b(t-\ell)
O G_{t,t-\ell}^{\partial}(\widehat{\mathbf w}_{t-\ell:t};\psi)
\mathbf s_k,
\label{eq:constructed_response}
\end{equation}
and stacking over time and source-major, basis-minor columns yields
\(H_\Phi^{\mathrm{lag}}\in\mathbb R^{N\times J}\), \(J=KB\), after applying
the observation selection from Equation~\ref{eq:observation_selection}.
Background correction gives
\begin{equation}
\widetilde Y=P_Q^\perp Y,
\qquad
\widetilde H_\Phi=P_Q^\perp H_{\Phi}^{\mathrm{lag}}.
\label{eq:constructed_projected_response}
\end{equation}

Scientific knowledge may declare a set \(\mathcal F_0\) of coefficients to be
identically zero before the response is diagnosed or fitted.  With
\(\mathcal I=\{1,\ldots,J\}\setminus\mathcal F_0\), IASA removes the
corresponding columns, computes every diagnostic from
\(\widetilde H_{\Phi,\mathcal I}\), fits only \(\mathbf c_{\mathcal I}\), and
then restores zeros in the full source--basis index.  The default is
\(\mathcal F_0=\varnothing\).  A small fitted coefficient is not a reason to
change this predeclared mask or recompute a more favorable diagnostic.
Source--basis coefficients are then estimated by nonnegative regularized least
squares,
\begin{equation}
\widehat{\mathbf c}
=\arg\min_{\mathbf c\ge 0}
\|\widetilde Y-\widetilde H_\Phi\mathbf c\|_2^2
+\lambda R(\mathbf c).
\label{eq:nnls_estimator}
\end{equation}
Reshaping \(\widehat{\mathbf c}\) into \(\widehat C\) reconstructs the
source activity trajectory
\begin{equation}
\widehat\theta_k(t)=\sum_{b=1}^{B}\widehat c_{kb}\phi_b(t).
\label{eq:reconstructed_activity}
\end{equation}
The regularizer can be ridge stabilization,
\begin{equation}
R(\mathbf c)=\|\mathbf c\|_2^2,
\label{eq:ridge_regularizer}
\end{equation}
or a weak prior around inventory-derived activity levels,
\begin{equation}
R(\mathbf c)=\|\mathbf c-\mathbf c_0\|_2^2.
\label{eq:prior_regularizer}
\end{equation}
Equivalently, one may fit background and source coefficients jointly:
\begin{equation}
(\widehat{\mathbf c},\widehat{\boldsymbol\gamma})
=\arg\min_{\mathbf c\ge 0,\boldsymbol\gamma}
\|Y-H_{\Phi}^{\mathrm{lag}}\mathbf c-Q\boldsymbol\gamma\|_2^2
+\lambda R(\mathbf c).
\label{eq:joint_fit}
\end{equation}
The projected formulation is used for identifiability because it exposes the
part of each source--basis fingerprint that remains after background correction.

\subsection{Constrained End-to-End Refinement as an Extension}
\label{subsec:constrained_refinement}

The reported estimator fixes wind fields and transport parameters before
fitting source activity.  A future PyTorch implementation may refine wind and
dispersion under explicit constraints.  Let \(\phi_0\) and \(\psi_0\) denote
the pretrained wind and initial dispersion parameters.  The prospective
refinement objective is
\begin{align}
\mathcal L_{\mathrm{refine}}
=&\|Y-H_{\Phi}^{\mathrm{lag}}(\phi,\psi)\mathbf c-Q\boldsymbol\gamma\|_2^2
+\lambda_\theta R(\mathbf c) \nonumber\\
&+\lambda_w\|\widehat{\mathbf w}_{1:T}^{\phi}-\widehat{\mathbf w}_{1:T}^{\phi_0}\|_2^2
+\lambda_\psi\|\psi-\psi_0\|_2^2 \nonumber\\
&+\lambda_{\mathrm{sm}}R_{\mathrm{sm}}(\widehat{\mathbf w}^{\phi}),
\label{eq:refinement_objective}
\end{align}
subject to
\begin{equation}
\mathbf c\ge 0,
\quad
\psi\in\Psi_{\mathrm{phys}},
\quad
\|\widehat{\mathbf w}_{1:T}^{\phi}-\widehat{\mathbf w}_{1:T}^{\phi_0}\|_\infty\le \epsilon_w.
\label{eq:refinement_constraints}
\end{equation}
This refinement is accepted only as a constrained local correction.  Without
these restrictions, wind and plume parameters could fit sparse pollution sensors
while weakening the physical and source-level interpretation of
\(\widehat{\mathbf c}\) and the reconstructed activities.
